
\begin{proposition}\label{cdp:st:D.1.12}
For a proper connected-fibre holomorphic map $f:Y\to C$ from a
compact complex threefold with $b_2(Y)=0$ and general fibre a
two-torus or a primary Kodaira surface, the cup-product homomorphism
is an isomorphism
\[
\bigwedge^2R^1f_*\OO_Y\xrightarrow{\simeq}R^2f_*\OO_Y.
\]
Every full fibre has structure-sheaf cohomology dimensions
$(h^0,h^1,h^2)=(1,2,1)$, and these two direct images satisfy
base change at every fibre. This proves these two cases of CDP20
Proposition 5.1 by an index and cup-product argument.
\end{proposition}

The torus-fibred six-sphere case and the cup-product index method are \citet[Theorem~2.3, pp.~5--6]{Viaclovsky2026Combs}. The statement here has the displayed broader hypotheses: $b_2=0$, an arbitrary base curve, and also primary Kodaira general fibres. The complete proof specifies the additional integral comparison for those hypotheses.

\begin{proof}
Write $E=R^1f_*\OO_Y$, which is locally free of rank two by \ref{cdp:st:E4},
and $G=R^2f_*\OO_Y$, which has rank one. Put
$d=\deg E$, $g=\deg(G/G_{\rm tor})$, and
$\ell=\operatorname{length}(G_{\rm tor})$.
Since $c_1(Y)=0$ rationally, Hirzebruch--Riemann--Roch gives
$\chi(Y,\OO_Y)=\int_Yc_1c_2/24=0$.
Coherent Leray and Riemann--Roch on $C$ give, with the rank terms
cancelled,
\[
0=\chi(\OO_C)-\chi(E)+\chi(G)=g+\ell-d.
\]
The cup product $\det E\to G$ is an isomorphism on a general
fibre. For a two-torus this is its exterior-algebra description.
For a primary Kodaira surface, use invariant $(1,0)$-forms
$\alpha,\eta$ with $d\alpha=0$ and
$d\eta=c\alpha\wedge\overline\alpha$, $c\ne0$.
Their conjugates are $\overline\partial$-closed. The class
$[\overline\alpha\wedge\overline\eta]$ is nonzero because its
Serre pairing with the holomorphic two-form $\alpha\wedge\eta$
is the nonzero integral of a volume form. This also proves that
$[\overline\alpha],[\overline\eta]$ are independent; since
$h^1(\OO)=2$ and $h^2(\OO)=1$, it proves the required cup-product
isomorphism. The invariant forms and their differential are
the forms in \cite[\S3, equations (9)--(12)]{Borcea1984}.
The composite of $\det E\to G$ with the torsion-free quotient is a
nonzero homomorphism of line bundles, so $g\geq d$. The displayed
identity and $\ell\geq0$ force $g=d$, $\ell=0$. A nonzero map
of line bundles of equal degree has no zero divisor, proving the
claimed global isomorphism.

Top-degree base change gives
$G\otimes\C(c)\simeq H^2(Y_c,\OO_{Y_c})$, of dimension one.
The multiplication-parameter exact sequence then identifies
$E\otimes\C(c)$ with $H^1(Y_c,\OO_{Y_c})$, since $G$ has no
parameter-torsion. \ref{cdp:st:E4} gives $h^0=1$. These arguments retain the
scheme-theoretic fibres throughout.
\end{proof}


